% Source: https://artofproblemsolving.com/community/c3h1236490_math_word_problem
% Stable identity: aops-topic-1236490-q01
The digits $2,3,4$ are used, so $A,B,C,D,E,F$ are $1,5,6,7,8,9$ in some order. Units: $E+F$ ends in $4$, so $E+F=14$ (the pairs from the available digits: $5+9$, $6+8$), with a carry of $1$. Tens: $C+3+1=D$ or $D+10$. Hundreds: $2+A+(\text{carry})=B$.

If $C+4=D$ (no carry): from the available digits, $(C,D)=(1,5)$, $(5,9)$; then $B=A+2$. With $(C,D)=(1,5)$ and $\{E,F\}=\{6,8\}$, the remaining digits $7,9$ give $A=7$, $B=9$, e.g. $216+738=954$: the digits $2,1,6,7,3,8,9,5,4$ are all different, so this works.

A short check of the remaining cases (with a carry in the tens column, or $(C,D)=(5,9)$) finds no other assignment. So $\overline{BD4}=\boxed{954}$.
